\documentclass[11pt]{article}
\usepackage[T1]{fontenc}
\usepackage{lmodern}
\usepackage[left=1in,right=1in,top=.84in,bottom=.88in]{geometry}
\usepackage{amsmath,amssymb,amsthm,booktabs,microtype}
\usepackage[hidelinks]{hyperref}
\newtheorem{theorem}{Theorem}
\newtheorem{lemma}[theorem]{Lemma}
\newtheorem{corollary}[theorem]{Corollary}
\newcommand{\N}{\mathbb N}
\DeclareMathOperator{\pos}{pos}
\title{Perfect mixing in increasing throwback queues}
\author{Blake Stone mathematics repository\\[-2pt]\small Proof developed with OpenAI Codex}
\date{November 26, 2026}
\begin{document}
\maketitle
\begin{abstract}
A throwback queue moves its first entry $x$ to position $x$, counting from zero.
For every strictly increasing sequence of positive integers we give the exact
return time of each finite configuration. The first $n+1$ labels visit every
admissible configuration precisely when $\gcd(x_j-j+1,x_i-i)=1$ for all
$0\le i<j\le n$. This answers the increasing-sequence perfect-mixing question
posed by Darnell and Dribus. For consecutive labels we determine all occurrence
times, the least period of the successive gaps, their reflection symmetry, and
the exact connection between first visits and Layman's sieve.
\end{abstract}

\section{Which queues mix perfectly?}
Let $x_0<x_1<\cdots$ be positive integers. Start with these labels in order.
At each integer time record the head $x$ and move it past the next $x$ entries.
A configuration of $x_0,\ldots,x_n$ records their absolute positions and
treats all other labels as indistinguishable. Each $x_i$ always has position at
most $x_i$: it starts at $i$, goes to $x_i$ on its own move, and otherwise
stays put or moves left. An admissible configuration assigns distinct positions
from $0$ to $x_i$ to these labels. Darnell and Dribus proved perfect mixing
for consecutive starting values in their Theorem 9, and asked which other
increasing sequences visit all their admissible configurations
\cite[pp.\ 15--16]{dd}.

Write $r_i=x_i-i+1$, so $r_i\ge2$. The empty prefix has return time
$P_{-1}=1$. Set $H_{-1}=1$. For $i\ge0$ put
\begin{equation}
 g_i=\gcd(r_i,H_{i-1}),\qquad
 P_i=P_{i-1}\frac{r_i}{g_i},\qquad
 H_i=\frac{H_{i-1}}{g_i}(r_i-1).
 \label{eq:period}
\end{equation}
\begin{theorem}\label{thm:main}
The first return of the configuration of $x_0,\ldots,x_n$ occurs at time
$P_n$. The orbit contains exactly $P_n$ configurations. It contains all
$\prod_{i=0}^{n}r_i$ admissible configurations if and only if
\begin{equation}
 \gcd(r_j,r_i-1)=1\qquad(0\le i<j\le n).
 \label{eq:mix}
\end{equation}
An infinite increasing queue mixes perfectly at every finite level exactly when
\eqref{eq:mix} holds for every pair $i<j$.
\end{theorem}

\begin{proof}
Erase the labels smaller than $x_i$ and let $z_i$ be the position of $x_i$
in the remaining queue. To reconstruct a configuration, place the labels in
increasing order. While placing $x_i$, exactly $i$ of the first $x_i+1$ slots
are occupied, because every earlier label has a smaller position bound. Thus
$z_i$ selects one of $r_i$ available slots, in order from left to right. This
gives a bijection with $\prod_{i=0}^{n}\{0,\ldots,r_i-1\}$. Initially all
coordinates are zero.

Consider erasing just the smallest label $s$. When $s$ leads, the relative
order of all remaining labels is unchanged. When a larger label $y$ leads,
it passes $s$: the old position of $s$ is between $1$ and $s<y$.
Erasing $s$ therefore reduces the jump of $y$ by exactly one. Subtracting
one from each remaining label gives a throwback queue again. The same
argument applies successively after erasing each smaller label.

Call a step active at level $i$ when no smaller label leads. At such a step
the coordinate $z_i$ decreases by one modulo $r_i$. Indeed, if $z_i=0$,
the reduced head is reinserted at $x_i-i=r_i-1$. If $z_i>0$, the reduced
leader passes it, moving it one slot left. At other steps $z_i$ stays fixed.
A step is active at level $i+1$ when it is active at $i$ and $z_i\ne0$.

Inductively, suppose the smaller coordinates form a cycle of length $P_{i-1}$
with $H_{i-1}$ steps active at $i$ per cycle. The assertion for the empty
prefix holds by convention. A return of the extended configuration must
first return the smaller coordinates. After $k$ of their cycles,
$z_i\equiv-kH_{i-1}\pmod{r_i}$. The least positive $k$ giving zero is
$r_i/g_i$. In those $k$ cycles there are $r_i H_{i-1}/g_i$ active steps.
Exactly one in each $r_i$ of them has $z_i=0$, so $H_i$ steps pass to the
next level. This proves both recurrences and the first return from time zero.
The update is deterministic, hence all states before this return are distinct.

At level $i$ a period can increase by at most $r_i$. A loss at any level
therefore precludes full capacity at every later level. If every earlier
level is full, its divisors $g_j$ are all one, and
$H_{i-1}=\prod_{j<i}(r_j-1)$. Full capacity extends one level precisely
when $r_i$ is coprime to this product. This is equivalent to
\eqref{eq:mix}.
\end{proof}

For $(1,3,4)$ the radices are $(2,3,3)$ and \eqref{eq:period} gives period
$18$, achieving the upper bound discussed in \cite[p.\ 16]{dd}. The prime
prefix $(2,3,5)$ has radices $(3,3,4)$. After the first two labels there
are four active steps in a period of nine. The third coordinate already
returns after this period, so exactly nine of the 36 configurations occur.
For a nonconsecutive infinite example, take $x_i=2^{i+1}+i-1$. All radices
are powers of two and every earlier $r_i-1$ is odd.

\begin{corollary}\label{cor:mean}
In a perfectly mixing prefix, the limiting mean absolute positions $M_i$
are determined successively by
\[
  M_i=\frac{x_i(x_i+1)/2-\sum_{j<i}M_j}{r_i}.
\]
\end{corollary}
\begin{proof}
A full period visits each admissible coordinate vector once. Conditional on
the placements of the smaller labels, $x_i$ is uniform on the $r_i$ remaining
slots in $\{0,\ldots,x_i\}$. Their sum is the sum of the whole interval
less the positions of the smaller labels. Take expectations.
\end{proof}

\section{Exact occurrence times for consecutive labels}
Suppose $x_j=m+j$ with $m\ge1$. Set $h=m+1$ and
\[
 Q(t)=t+1+\lfloor t/m\rfloor\qquad(t\in\mathbb Z).
\]
On nonnegative integers, $Q$ lists in increasing order the times not
divisible by $h$, with $Q(0)=1$. In fact $Q(ma+b)=ha+b+1$ for
$0\le b<m$. The smallest label $m$ leads exactly at times $hk$.
After erasing this label and subtracting one, the initial queue is again
$m,m+1,\ldots$. The erasure argument in Theorem~\ref{thm:main}
therefore identifies the leader at time $Q(t)$ as one more than the leader
at time $t$. This is also the consecutive-label instance of the placement
construction of Darnell and Dribus \cite[Sec.\ 4]{dd}.

\begin{theorem}\label{thm:gaps}
Let $E_j(k)$ be the time of occurrence number $k+1$ of label $m+j$, and put
$G_j(k)=E_j(k+1)-E_j(k)$, with $k\ge0$. Then
\[
 E_j(k)=Q^{\circ j}(hk).
\]
The gap sequence has least positive period $D=m^j$. Its first $D-1$ gaps
form a palindrome. Its last gap in the first period is $h+2E_j(0)$.
\end{theorem}
\begin{proof}
The occurrence formula follows by induction from the description of $Q$.
Extend the formula to integer $k$. Euclidean division gives
\[
 Q(t+mz)=Q(t)+hz,\qquad Q(-h-t)=-h-Q(t).
\]
For the reflection, write $t=ma+b$ with $0\le b<m$. The remainder of
$-h-t$ is $m-1-b$. Iterating the translation gives
$Q^{\circ j}(t+m^jz)=Q^{\circ j}(t)+h^jz$, and the reflection commutes
with every iterate. Consequently
\[
 E_j(k+D)=E_j(k)+h^{j+1},\qquad E_j(-1-k)=-h-E_j(k).
\]
The first identity yields a gap period $D$ of total weight $h^{j+1}$.
For any positive period $p$, summing the gaps over $pD$ terms in periods
of both lengths gives $D\mid p h^{j+1}$. Since $\gcd(D,h)=1$, $D\mid p$.
The two identities also give $G_j(D-2-k)=G_j(k)$ and
$G_j(D-1)=G_j(-1)=h+2E_j(0)$. These verify the assertions including $D=1$.
\end{proof}

Darnell and Dribus identified reflection about $-2$ for $m=3$ in their two-sided placement construction \cite[Sec.\ 4]{dd}. The formula above gives its exact action on the gap indices for every $m$.
For $m=3$ this proves Kozar's period and palindrome observations for
\href{https://oeis.org/A357081}{A357081} \cite{kozar}. For example, the
gaps between appearances of $5$ repeat as
$7,6,8,7,7,8,6,7,8$. If $L(t)$ is the leader, then $L(4a)=3$ and
$L(t)=1+L(t-\lceil t/4\rceil)$ for $4\nmid t$. After shifting the
index by one this is precisely the defining recurrence for
\href{https://oeis.org/A087165}{A087165}; hence $L(t)=2+\mathrm{A087165}(t+1)$.
Removing all threes raises the remaining sequence by one relative to the
original, and applying the same deletion rule twice to A087165 shows that
removing its ones and twos gives A357081. For $m=1$ the formula recovers
the ruler sequence: $E_j(k)=2^{j+1}k+2^j-1$.

\section{First visits and the sieve identification}
The first visit to $m+j$ is at $F_j=Q^{\circ j}(0)$. These times increase
strictly, so they are exactly the record times. In particular
\[
 F_0=0,\qquad F_{j+1}=1+\left\lfloor\frac{(m+1)F_j}{m}\right\rfloor.
\]
To compare with Layman's $L$-sieve, remove the first item of the current
list, then delete positions $h-1,2h-1,\ldots$ in the remaining list,
reindexing from one after removing its head. The $j$-th survivor among
the tail positions has index $b(j)=\lfloor hj/(h-1)\rfloor$ for $j\ge1$:
in each block of $h$ one drops position $h-1$ and retains the final
position. Thus the next list, as a subsequence of the previous list,
has index map $j\mapsto1+b(j)=Q(j)$. Repeating that same deletion
gives successive first elements $1,Q(1),Q^{\circ2}(1),\ldots$.
Since $Q(0)=1$, these are $F_1,F_2,\ldots$.
For $m=3$ this proves the first-visit identification with
\href{https://oeis.org/A155167}{A155167} noted by Sloane and recorded in
Darnell and Dribus's Table 3. The initial record at time zero is separate.
Infinite recurrence of the labels was proved for general distinct positive
inputs in \cite[Thm.\ 2]{dd}.

\begin{thebibliography}{9}
\bibitem{dd} J. Darnell and B. F. Dribus, \emph{Throwback Sequences of
Positive Integers}, Journal of Integer Sequences 28 (2025), Article 25.5.1.
\href{https://cs.uwaterloo.ca/journals/JIS/VOL28/Dribus/dribus4.pdf}{Published paper}.
\bibitem{kozar} A. M. Kozar Jr., OEIS \href{https://oeis.org/A357081}{A357081}
(2022). See also \href{https://oeis.org/A087165}{A087165} for the
earlier recurrence and \href{https://oeis.org/A155167}{A155167} for the sieve.
\end{thebibliography}
\end{document}
